\documentclass[11pt]{article}
\usepackage{amsmath,amssymb,amsthm}
\newtheorem{proposition}{Proposition}
\begin{document}
\title{Scope limits of some auxiliary-spin approaches to AC2 upper bounds}
\author{}\date{}\maketitle

This note records negative results about specified proof methods. It does
not prove that the current unrestricted AC2 upper bound is optimal, and
it does not prove a stronger unrestricted upper bound.

For an AC2 code of spin $j$, coupling to an auxiliary spin $s$ makes
\[
 W=X^2+aX-\frac{j(j+1)s(s+1)}3I,\qquad X=\mathbf J\cdot\mathbf S,
\]
have zero compression on the code tensored with the auxiliary space.
The eigenvalue in total spin $L$ is obtained by substituting
\[
 x_L=\frac{L(L+1)-j(j+1)-s(s+1)}2.
\]
The exact script \texttt{ancilla\_quadratic\_scan.py} checks every
nonsingular sign pattern as $a$ varies, for $8\le N\le100$ and
$s=1/2,1,\ldots,20$. No bound in this scan improves
\[
 \min\left\{\frac{N-2}{4},\frac{N+5}{5}\right\}.
\]
Only finitely many ancillary spins and physical sizes are covered by that
statement; it is not a theorem of finite-size optimality for this method.

\begin{proposition}
For any fixed irreducible ancillary spin $s>2$, the quadratic witness
family above cannot improve the leading coefficient $1/5$ of the
spin-two bound as $j\to\infty$.
\end{proposition}
\begin{proof}
In the branches $L=j+m$, $m=-s,-s+1,\ldots,s$, we have
$x_L=jm+O(1)$. If $a/j$ has a finite limiting value $A$, the limiting
quadratic is
\[
 p(m)=m^2+Am-\frac{s(s+1)}3.
\]
Its roots are $-u<0<v$, with $uv=s(s+1)/3$. Allowing either sign at a
zero of the limiting quadratic only makes the following counting bounds
weaker, and therefore valid for finite-size perturbations.

Let $t$ grid points lie strictly outside the closed root interval. If
there are $a_-$ removed points at the left end and $a_+$ at the right,
and both surviving endpoints straddle zero, inclusion of those endpoints
requires
\[
 (s-a_-)(s-a_+)\le\frac{s(s+1)}3.
\]
Since the product is at least $s^2-s(a_-+a_+)$, this implies
\[
 t=a_-+a_+\ge\frac{2s-1}{3}>\frac{2s+1}{5}\qquad(s>2).
\]
If a surviving endpoint does not straddle zero, at least half the grid
has already been removed, which gives the same weaker conclusion.

The open root interval intersected with $[-s,s]$ has length at least
$s$: if neither root is clipped this follows from
$u+v\ge2\sqrt{s(s+1)/3}>s$; if one root is clipped, the interval still
contains a full half of $[-s,s]$. Hence at least
$\lceil s\rceil-1\ge s-1>(2s+1)/5$ grid points lie strictly inside it.
Both sign counts therefore exceed $(2s+1)/5$. Each branch has
multiplicity $N+O(1)$, so the inertia bound divided by the auxiliary
dimension $2s+1$ has leading coefficient strictly greater than $1/5$.
If $|a|/j\to\infty$, the linear term determines the signs away from
$m=0$ and gives approximately half the branches of either sign, also
larger than $1/5$. Subsequences cover varying choices of $a$.
\end{proof}

\section*{An obstruction for one promising reducible pair}
The preceding proposition says nothing about reducible ancillary spaces
or non-scalar matrix coefficients. One candidate is $V_3\oplus V_4$, of
dimension sixteen, for which three positive branches could give a
leading coefficient $3/16<1/5$.

For the natural one-sided sign pattern, the two classical diagonal blocks
can be normalized to
\[
 A_m=m^2+am-4,\quad B_m=m^2+bm-20/3,
 \quad 5/3\le a\le3,\quad7/3\le b\le17/3.
\]
Both are nonpositive at $m=-3,\ldots,1$ and positive at $m=2,3$; $B_4$ is
also positive. A general rank-one/rank-two off-diagonal coupling has the
form $\sqrt{16-m^2}\,\ell_m$, where $\ell_m=e+fm$ and $e,f$ may be
complex. Avoiding new positive eigenvalues in the blocks $m=\pm1$ requires
\[
 |\ell_1|^2\le\frac{(3-a)(17/3-b)}{15},\quad
 |\ell_{-1}|^2\le\frac{(3+a)(17/3+b)}{15}.
\]
Using $\ell_3=2\ell_1-\ell_{-1}$ and Cauchy--Schwarz gives
\[
 |\ell_3|^2\le\frac{34}{15}(5-a)\le\frac{68}{9}.
\]
But merging the two positive eigenvalues of the $m=3$ block into a single
positive eigenvalue requires its determinant to be negative, hence
\[
 |\ell_3|^2>
 \frac{(5+3a)(7/3+3b)}7\ge\frac{40}{3}>\frac{68}{9}.
\]
This contradiction rules out that particular three-positive-branch
pattern. It does not exclude all reducible ancillary constructions.
\end{document}
